\section{Syntax of ILD}
\label{syntax}
% PAPER-REVIEW[PR-STYLE-0003] pass=style state=open
% Issue: The same set Val is called "the internal syntax", "tangible values" (this sentence) and "Observable values" (Grammar 2.1). The three terms suggest three different notions.
% Evidence: "a subset of the internal syntax (tangible values) \ildsetVal" vs. grammar annotation "Observable values".
% Proposal: Choose one term and use it in both places, e.g. "a subset of the set of values \ildsetVal (Grammar 2.1)", with "Values" as the grammar annotation.
% Reason: Consistent terminology for the central semantic domain.
% END-PAPER-REVIEW[PR-STYLE-0003]
We define ILD as a homoiconic language where parseable programs (\(\ildsetProg\))
are a subset of the internal syntax (tangible values) \(\ildsetVal\).
% PAPER-REVIEW[PR-TECHNICAL-0005] pass=technical state=open
% Issue (demonstrated, rendering): The footnote on "Binding environments" in Grammar 2.1 is never typeset, so its assumption (unique keys) and the remark on other representations are missing from the PDF.
% Evidence: The PDF from latexmk -lualatex (2026-11-26) shows the superscript marker, but pdftotext finds no "We assume keys are restricted to be unique". Standard tabular, here inside cfighere/center, drops \footnote text. Eq. (1) ("e = (... (s . v) ...)") and eq. (9) rely on this assumption.
% Proposal: Use \footnotemark inside the tabular and \footnotetext{...} after \end{cfighere}, or move the two sentences into the prose after the grammar. Also say that eq. (9) preserves uniqueness, which is what makes the assumption hold for environments built by application.
% Reason: An assumption the semantics depends on is currently invisible to readers. Check the vtu2docx output as well.
% END-PAPER-REVIEW[PR-TECHNICAL-0005]
\begin{cfighere}{Grammar}{grammar}{Syntax of ILD}{BNF definition of the ILD syntax}\label{syntax-grammar}
{\footnotesize
\begin{tabular}{r@{\;}c@{\;}l@{\qquad}l}
\(\ildsetVal\) & \(:=\) & \(\ildsetSexp \mid \ildsetSF \mid \ildsetFail \mid \ildsetFunc \mid\) & Observable values \\
 & & \(\ildsetSym \mid \ildsetNum \mid \ildsetBool \mid \ildsetStr\) & \\
\(\ildsetSym\) & \(\in\) & \emph{symbols} & \\
\(\ildsetSexp\) & \(:=\) & \(() \mid \ildpair{\ildsetVal}{\ildsetVal}\) & S-expressions \\
\(\ildsetSF\) & \(:=\) & \(\ildsf{free-vars} \mid \ildsf{quote} \mid \ildsf{macroexpand}\) & Special forms \\
\(\ildsetFail\) & \(:=\) & \(\ildfail{\ildsetVal}\) & Failure objects, each carrying a context value \\
\(\ildsetFunc\) & \(:=\) & \(\ildsetBuiltinFunc \mid \ildsetUserFunc\) & Ordinary functions \\
\(\ildsetBuiltinFunc\) & \(:=\) & \(\ildhost{+} \mid \ildhost{cons} \mid \ildhost{apply} \mid \ildhost{eval} \mid \dots\) & Builtin functions, \cref{builtin-funcs} \\
\(\ildsetUserFunc\) & \(:=\) & \(\ildabstr{\ildsetEnv}{\ildsetSymList}{\ildsetVal}\) & User-defined functions, \cref{abstraction} \\
\(\ildsetNum\) & \(\in\) & \emph{numbers} & \\
\(\ildsetBool\) & \(:=\) & \(\ildhost{t} \mid \ildhost{f}\) & \\
\(\ildsetStr\) & \(\in\) & \emph{strings} & \\
\(\ildsetProg\) & \(:=\) & \(\ildsetSym \mid () \mid \ildpair{\ildsetProg}{\ildsetProg} \mid\) & Surface syntax (programs) \\
 & & \(\ildsetStr \mid \ildsetNum \mid \ildsetBool\) & \\
\(\ildsetList\) & \(:=\) & \(() \mid \ildpair{\ildsetVal}{\ildsetList}\) & S-expression \emph{lists} \\
\(\ildsetEnv\) & \(:=\) & \(() \mid \ildpair{\ildpair{\ildsetSym}{\ildsetVal}}{\ildsetEnv}\) & Binding environments\footnote{We assume keys are restricted to be unique (\cref{semantics}). Furthermore, performant implementations will use other representations of binding environments.} (K/V lists) \\
\(\ildsetSymList\) & \(:=\) & \(() \mid \ildpair{\ildsetSym}{\ildsetSymList}\) & Symbol lists (parameters, argument lists) \\
\end{tabular}}
\end{cfighere}
% PAPER-REVIEW[PR-LANGUAGE-0003] pass=language state=open
% Issue: "i.e." and "e.g." followed by a normal space are typeset with inter-sentence spacing (visible gap in the PDF, e.g. "(e.g.  Listing 4.5)").
% Evidence: "i.e. the nested pairs" (this paragraph); "constructed via e.g. \ild{lambda}" (§3.3 list); "(e.g. \cref{if})" (§3.3.1).
% Proposal: Write "i.e.,~"/"e.g.,~" or "e.g.\ ", or add \frenchspacing to the preamble. Check that vtu2docx handles the chosen form before applying.
% Reason: Typographic correctness of the PDF submission; not applied directly because the Word converter's handling of "\ " was not verified.
% END-PAPER-REVIEW[PR-LANGUAGE-0003]
Source code is parsed as
standard S-expressions. Throughout this paper, we use the \(\ildlist{a_1}{a_2}{\dots}{a_n}\)
syntax to denote S-expression lists, i.e. the nested pairs
\(\ildpair{a_1}{\ildpair{a_2}{\dots \ildpair{a_n}{()} \dots}}\)~\cite{sexp}.
In addition, the parser of ILD is augmented with the following syntactic shorthands:
\begin{itemize}
  \item Quote: \ild{'<expr>} \(\Mapsto\) \ild{(quote <expr>)}
  \item Macroexpand: \\
    \ild{(!<expr1> ... <exprN>)} \(\Mapsto\) \\
    \ild{(macroexpand <expr1> ... <exprN>)}
\end{itemize}
\section{Semantics of ILD}
\label{semantics}
We define the semantics of ILD expressions via the \(\denot{}{\cdot}\) notation:
\(\denot{e}{v}\) stands for “evaluating the value \(v\) in the binding environment
\(e\)”.

Symbols (\( s \in \ildsetSym \)) are looked up in the binding environment:
\begin{equation}
  \denot{e}{s} = \begin{cases}
      v & e = \ildlist{\dots}{\ildpair{s}{v}}{\dots} \\
      \ildfail{\text{unbound symbol}} & e \neq \ildlist{\dots}{\ildpair{s}{v}}{\dots} \\
    \end{cases}
\end{equation}
% PAPER-REVIEW[PR-TECHNICAL-0004] pass=technical state=open
% Issue (minor inconsistency): The prose says every value that is neither a symbol nor a pair evaluates to itself. The empty list () is neither, yet §3.4 says it evaluates to #fail. The set in parentheses correctly excludes ().
% Evidence: "Values that are neither symbols nor pairs (...) evaluate to themselves" vs. §3.4 "This includes the empty list".
% Proposal: "Values other than symbols, pairs and the empty list (...) evaluate to themselves", or "Atoms other than symbols and () ...".
% Reason: The prose and the formula should define the same case split.
% END-PAPER-REVIEW[PR-TECHNICAL-0004]
Values that are neither symbols nor pairs (\(v \in \ildsetSF \cup \ildsetFail \cup \ildsetNum \cup \ildsetBool \cup \ildsetStr \cup \ildsetFunc\))
evaluate to themselves:
\begin{equation}
  \denot{e}{v} = v
\end{equation}
Proper lists (also known as \emph{combinations}) are evaluated according to the following rules:
\begin{equation}
  \begin{gathered}
  \denot{e}{\ildlist{h}{\alpha_1}{\dots}{\alpha_n}}
  = \begin{cases}
      \denotsf{e}{\denot{e}{h}}{\alpha_1}{\dots}{\alpha_n} & \denot{e}{h} \in \ildsetSF \\
      \denotapply{\denot{e}{h}}{\denot{e}{\alpha_1}}{\dots}{\denot{e}{\alpha_n}} & \denot{e}{h} \in \ildsetFunc \\
      \ildfail{\text{cannot apply }\denot{e}{h}} & \text{otherwise} \\
    \end{cases}
  \end{gathered}
\end{equation}
Depending on whether the combination head evaluates to a special form (\(\ildsetSF\))
or an ordinary function (\(\ildsetFunc = \ildsetBuiltinFunc \cup \ildsetUserFunc \)),
it is \emph{applied} via either \(\interopword{sf}\) or \(\interopword{apply}\) respectively.
Ordinary function application and the special forms are outlined in the following sections.
In order for ILD programs to be able to access special forms (\(\ildsetSF\)) and builtin functions
(\(\ildsetBuiltinFunc\)), symbols that resolve to them are provided in the root environment \(E\):
\begin{equation}
  \begin{array}{@{}l@{}}
    E = ( \\
    \hspace{3em} \ildpair{\ildsym{quote}}{\ildsf{quote}} \\
    \hspace{3em} \ildpair{\ildsym{macroexpand}}{\ildsf{macroexpand}} \\
    \hspace{3em} \ildpair{\ildsym{free-vars}}{\ildsf{free-vars}} \\[1em]
    \hspace{3em} \ildpair{\ildsym{cons}}{\ildhost{cons}} \\
    \hspace{3em} \ildpair{\ildsym{car}}{\ildhost{car}} \\
    \hspace{3em} \ildpair{\ildsym{cdr}}{\ildhost{cdr}} \\
    \hspace{3em} \ildpair{\ildsym{+}}{\ildhost{+}} \\
    \hspace{3em} \ildpair{\ildsym{mk-lambda}}{\ildhost{mk-lambda}} \\
    \hspace{3em} \ildpair{\ildsym{make-fail}}{\ildhost{make-fail}} \\
    \hspace{3em} \ildpair{\ildsym{read-source}}{\ildhost{read-source}} \\
    \hspace{3em} \dots \\
    \hspace{1em})
  \end{array}
\end{equation}
\subsection{Special forms}
% PAPER-REVIEW[PR-STYLE-0002] pass=style state=open
% Issue: "Unlike Scheme or Common LISP, ILD is a simple language with only three special forms" implies that Scheme and Common Lisp are not simple languages, which is probably not the intended contrast.
% Evidence: The contrast intended is the number of special forms.
% Proposal: "Unlike Scheme or Common LISP, which define many special forms, ILD has only three."
% Reason: The comparison should be limited to what is actually being compared.
% END-PAPER-REVIEW[PR-STYLE-0002]
Applying a special form is done on the \textbf{unevaluated} arguments of the combination.
Unlike Scheme or Common LISP, ILD is a simple language with only three special forms.
\subsubsection{The quote special form}
The \(\ildsf{quote}\) special form returns its argument without evaluating it.
\begin{equation}
    \denotsf{e}{\ildsf{quote}}{\alpha} = \alpha
\end{equation}
In real-world code, the \ild{'<expr>} syntax is often used instead of
\ild{(quote <expr>)} (see \cref{syntax}).
\subsubsection{The free-vars special form}
The \(\ildsf{free-vars}\) special form captures the current binding environment and is used by higher-level constructs
such as the lambda macro (\cref{lambda-macro}).
\begin{equation}
    \denotsf{e}{\ildsf{free-vars}} = e
\end{equation}
\subsubsection{Macro expansion}
\label{macroexpand-mechanism}
% PAPER-REVIEW[PR-TECHNICAL-0007] pass=technical state=open
% Issue (claim not demonstrated; question for the author): The paper's central claim, that ILD keeps static scoping in programs that use macros, is asserted but never argued, and the semantics seems to qualify it.
% Evidence: Claim: intro.tex ¶2 ("does not sacrifice expressiveness of metaprogramming while retaining immutability and static scoping"), the title, and the conclusion ("statically-scoped programs"). Mechanism: eq. (7) evaluates a macro's output in the caller's environment e, and this section compares the behaviour to unhygienic macro systems, so a symbol introduced by an expansion resolves in the caller's scope. The lambda and if macros avoid this only because they splice values (#mk-lambda, special-form values) instead of symbols (Listings 4.2 and 4.5), a discipline the paper never states. Environments are also first-class data: free-vars exposes e, and mk-lambda and eval accept any Env value.
% Proposal: (a) Define "statically scoped" for ILD precisely (e.g. every symbol in a function body resolves in the closure's environment extended with the parameters, eq. (8); there is no dynamic binding). (b) State value-splicing as the means by which macros avoid depending on the caller's bindings, and say what it does not guarantee (hygiene: a macro-introduced binder can still capture user identifiers). (c) Either argue that first-class environments (free-vars, mk-lambda, eval) do not break the property, or qualify the claim.
% Reason: This is what distinguishes ILD from Lispkit Lisp and from unhygienic Lisps; reviewers will look for the argument here. Related: PR-STRUCTURE-0004, PR-TECHNICAL-0010.
% END-PAPER-REVIEW[PR-TECHNICAL-0007]
Macros in ILD are plain functions (\(\varphi \in \ildsetFunc\)) that assume their arguments
are program fragments. We refer to such functions as \emph{macro functions}.
To invoke a function as a macro, the special form
\ildsf{macroexpand} is used at the callsite:
\begin{equation}
  \denotsf{e}{\ildsf{macroexpand}}{\varphi}{\vec{\alpha}}
  = \denot{e}{\denotapply{\denot{e}{\varphi}}{\vec{\alpha}}}
\end{equation}
An expression
\(\ildlist{\ildsf{macroexpand}}{\varphi}{\alpha_1}{\alpha_2}{\dots}{\alpha_n}\)
evaluates just \(\varphi\) and then passes the
\textbf{unevaluated} arguments to it. This allows the macro function to treat the subprogram
passed to it as data and transform it arbitrarily; the resulting program is then
evaluated in turn. Macro expansion, being a regular evaluation step,
proceeds from the outermost call inward, similarly to macro semantics
in most LISP dialects with unhygienic macro systems~\cite{steele-gabriel-evolution,kohlbecker1986syntactic}.
\subsection{User-defined functions}
\label{abstraction}
% PAPER-REVIEW[PR-TECHNICAL-0012] pass=technical state=open
% Issue (notation): \subst{x}{v} renders as x stacked over v. The widespread substitution notation [v/x] reads "v replaces x", so e[p_1 over a_1] in eq. (8) is easily read backwards ("p_1 replaces a_1"). The body is also named β in eq. (8) but B in §3.3 (mk-lambda).
% Evidence: Eqs. (8)-(9); §3.3 "apply(#mk-lambda; e, P, B)".
% Proposal: Use an unambiguous binding notation such as e[x ↦ v] (or e[v/x] with one sentence fixing the convention), and use one letter for the body.
% Reason: Avoids a likely misreading of the core application rule.
% END-PAPER-REVIEW[PR-TECHNICAL-0012]
User-introduced functional objects are
represented as \(\ildsetUserFunc\)-terms. Each such term carries a binding environment \(e\),
a list of formal parameters \(P \in \ildsetSymList\) and a body expression
\(\beta\), and is \emph{applied} by substituting the formal parameters by the
actual parameters in the local binding environment,
and then evaluating the body in the resulting environment:
\begin{equation}\begin{gathered}
  \denotapply{\ildabstr{e}{P = \ildlist{p_1}{p_2}{\dots}{p_n}}{\beta}}{a_1}{a_2}{\dots}{a_n} = \denot{e'}{\beta}
  \\
  \text{where } e' = \spacedlist{e}{\subst{p_1}{a_1}}{\subst{p_2}{a_2}}{\dots}{\subst{p_n}{a_n}}
\end{gathered}\end{equation}
Formally, parameter substitution in a binding environment \(e\subst{x}{v}\) is defined as:
\begin{equation}
  e \subst{x}{v} =
  \begin{cases}
    \ildpair{\ildpair{x}{v}}{()} & e = () \\
    \ildpair{\ildpair{x}{v}}{e'} & e = \ildpair{\ildpair{x}{w}}{e'} \\
    \ildpair{\ildpair{y}{w}}{e' \subst{x}{v}} & e = \ildpair{\ildpair{y}{w}}{e'} \land y \neq x
  \end{cases}
\end{equation}
Such functional objects are created via the \(\ildsym{mk-lambda}\) constructor (\cref{builtin-funcs})
and then used for \ild{lambda}-expressions (\cref{lambda-macro}).
\subsection{Builtin functions}
\label{builtin-funcs}
The ILD interpreter and root environment provide a set of builtin functions
in order for ILD programs to be able to operate on the builtin data structures
(\(\ildsetNum\), \(\ildsetBool\), \(\ildsetStr\)):
\begin{itemize}
  \item Arithmetic: \(\ildhost{+}\), \(\ildhost{*}\), \(\ildhost{-}\), \(\ildhost{/}\)
  \item Comparison: \(\ildhost{=}\), \(\ildhost{<}\),
    \(\ildhost{>}\), \(\ildhost{<=}\), \(\ildhost{>=}\)
  \item Pairs:
    \begin{itemize}
      \item \(\denotapply{\ildhost{cons}}{a}{d} = \ildpair{a}{d}\)
      \item \(\denotapply{\ildhost{car}}{\ildpair{a}{d}} = a\)
      \item \(\denotapply{\ildhost{cdr}}{\ildpair{a}{d}} = d\)
    \end{itemize}
  \item Constructor for lists: \(\ildhost{list}\)
  \item User-defined functions: \(\denotapply{\ildhost{mk-lambda}}{e}{P}{B} = \ildabstr{e}{P}{B}\) \\
  constructed via e.g. \ild{lambda}-expressions (\cref{lambda-macro})
  \item Value type predicates: \(\ildhost{null?}\), \(\ildhost{pair?}\),
    \(\ildhost{symbol?}\), \(\ildhost{string?}\), \(\ildhost{func?}\),
    \(\ildhost{fail?}\)
  \item Equality: \(\denotapply{\ildhost{sym-eq?}}{s_1}{s_2}\), \(\denotapply{\ildhost{str-eq?}}{s_1}{s_2}\)
  \item Failure: \(\denotapply{\ildhost{make-fail}}{v} = \ildfail{v}\)
\end{itemize}
\subsubsection{Branching}\label{branching}
The boolean datatype is of no use on its own without a function that
distinguishes between its values. Thus, we provide
the \ildhost{bool-to-k} function that transforms a boolean into
an argument selector:
\begin{equation}
  \denotapply{\ildhost{bool-to-k}}{b} =
      \begin{cases}
        \ildabstr{()}{\ildlist{x}{y}}{x} & b = \ildhost{t} \\
        \ildabstr{()}{\ildlist{x}{y}}{y} & b = \ildhost{f} \\
      \end{cases}
\end{equation}
The \ildhost{bool-to-k}
function can then be used to build higher-level branching primitives
(e.g. \cref{if}).
\subsubsection{A recursion operator}
% PAPER-REVIEW[PR-TECHNICAL-0003] pass=technical state=open
% Issue: The poly-fix definition has an index error and an unexplained step, and its notation collides with the boolean #f.
% Evidence: (a) "where f_1 \dots f_i are intermediate recursion results" should run to f_n. (b) "intermediate recursion results of the fixed point" does not say why eqs. (12)-(13) are well founded: each #f_i is an opaque host object that, when applied, unfolds Γ_i once with all #f_j as arguments. That is why the definition terminates under the call-by-value evaluation of eq. (3), provided each Γ_i returns a function. (c) \ildconst{f}{i} renders as "#f_1", while \ildhost{f} ("#f") is the false boolean in Grammar 2.1 and eq. (10).
% Proposal: Use f_1 \dots f_n. Replace the phrase with an explicit sentence along the lines of (b), and state the requirement that each Γ_i returns a function. Rename the host objects (e.g. #r_i or #fix_i).
% Reason: poly-fix is the only source of recursion in the language and the basis of Y and letrec.
% END-PAPER-REVIEW[PR-TECHNICAL-0003]
We provide a low-level operator \ildhost{poly-fix} for constructing
mutually-recursive functions:
\begin{equation}
  \denotapply{\ildhost{poly-fix}}{\Gamma_1}{\Gamma_2}{\dots}{\Gamma_n} = \ildlist{\ildconst{f}{1}}{\ildconst{f}{2}}{\dots}{\ildconst{f}{n}}
\end{equation}
When evaluated, \ildhost{poly-fix} introduces the functional objects \(\ildconst{f}{1}, \dots, \ildconst{f}{n}\) whose semantics are:
\begin{equation}
  \denotapply{\ildconst{f}{i}}{a_1}{\dots}{a_k}
  = \denotapply{f_i}{a_1}{\dots}{a_k}
\end{equation}
where \( f_1 \dots f_i \) are intermediate recursion results of the fixed point:
\begin{equation}
  f_i = \denotapply{\Gamma_i}{\ildconst{f}{1}}{\ildconst{f}{2}}{\dots}{\ildconst{f}{n}}
\end{equation}
\subsubsection{Interpreter access}
% PAPER-REVIEW[PR-TECHNICAL-0002] pass=technical state=open
% Issue (demonstrated inconsistency): The equation for #eval swaps expression and environment: apply(#eval; e, v) = \denot{v}{e} renders as ⟦e⟧_v, i.e. "evaluate e in environment v".
% Evidence: §3 defines ⟦v⟧_e as evaluating v in environment e (\denot{e}{v}). Listing 4.3 calls (eval (free-vars) (read-source ...)) with the environment first and the program second.
% Proposal: Write \denot{e}{v} on the right-hand side, i.e. ⟦v⟧_e. Optionally use the letters ε/φ or env/expr to avoid reusing v and e.
% Reason: eval is used in the central bootstrapping step (Listing 4.3).
% END-PAPER-REVIEW[PR-TECHNICAL-0002]
It is practical to allow programs to call into the interpreter:
\begin{equation}\begin{gathered}
  \denotapply{\ildhost{apply}}{f}{\ildlist{a_1}{\dots}{a_n}} = \denotapply{f}{a_1}{\dots}{a_n}\\
  \denotapply{\ildhost{eval}}{e}{v} = \denot{v}{e}
\end{gathered}\end{equation}
% PAPER-REVIEW[PR-TECHNICAL-0008] pass=technical state=open
% Issue (tension with the purity claim): read-source makes evaluation depend on file-system state, while the paper claims the language is purely functional and fully immutable, and lists IO as future work.
% Evidence: "(read-source p) will read the file whose path is p, parse its contents and return the resulting value" (this paragraph); "purely-functional", "fully immutable" (title, intro.tex ¶2, conclusion); "ILD can be extended with ... side effects, IO" (§5.1).
% Proposal: State the assumption that makes read-source pure, e.g. "read-source is a pure lookup in a fixed, immutable map from program names to programs; reading files is an implementation of that map, and we assume the source tree does not change during evaluation". Alternatively, name read-source explicitly as the only exception.
% Reason: Without the assumption, a reviewer can point to read-source as a counterexample to the paper's main property.
% END-PAPER-REVIEW[PR-TECHNICAL-0008]
Furthermore, to allow programs to reference other programs, we provide
a builtin function \ildhost{read-source} that retrieves a program by name.
In real-life implementations, \(\ildapp{read-source}{p}\) will read
the file whose path is \(p\), parse its contents and return the resulting value.
\subsection{Improper constructs}
% PAPER-REVIEW[PR-TECHNICAL-0006] pass=technical state=open
% Issue (missing specification): Failure behaviour is defined only for evaluation cases, not for application cases, and it is unclear whether failures propagate.
% Evidence: (a) Eq. (8) defines apply for a user function only when the number of arguments equals |P|; arity mismatch is not covered, and this section speaks only of "\ildsetVal-cases" and special forms. (b) Builtins with ill-typed arguments (car of a non-pair, + on strings, bool-to-k on a non-boolean) are unspecified. (c) Eq. (3) evaluates the arguments and applies the function even when an argument is a #fail value, so it is open whether e.g. (cons #fail⟨x⟩ 1) is a pair or a failure. Listing 4.6 returns make-fail as an ordinary value. (d) macroexpand whose head evaluates to a non-function (e.g. a special form) is not covered by eq. (7).
% Proposal: Add one rule covering apply, e.g. "every application not covered above, including arity mismatches and ill-typed builtin arguments, yields #fail⟨...⟩". State explicitly whether failures are ordinary first-class values or short-circuit evaluation. Confirm against the implementation.
% Reason: Completeness of the semantics is part of the paper's contribution; readers cannot otherwise predict the result of common errors.
% END-PAPER-REVIEW[PR-TECHNICAL-0006]
All \(\ildsetVal\)-cases that are not listed in the previous sections evaluate
to a \(\ildfailbare\). This includes the empty list, non-list pairs and special forms
applied to a wrong number of operands.
